\documentclass[a4paper,hidelinks,14pt]{extarticle}

\usepackage[utf8]{inputenc}
\usepackage[T2A]{fontenc}
\usepackage[english, russian]{babel}
\usepackage{lipsum}
\usepackage{amsmath}
\usepackage{amssymb}
\usepackage{amsfonts}
\usepackage{mathtools}
\usepackage{datetime}
\usepackage[pdftex]{graphicx}
\usepackage{indentfirst}
\usepackage{asymptote}
\usepackage{systeme}
\usepackage[dvipsnames]{xcolor}
\usepackage{lastpage}
\usepackage{fancybox,fancyhdr}
\usepackage{hyperref}
\usepackage[font={small,it}]{caption}
\fancyhead[L]{Лабораторная работа №6}
\fancyhead[C]{}
\fancyhead[R]{\textit{$H_\infty$-оптимизация}}
\fancyfoot[L]{}
\fancyfoot[C]{\thepage\space}
\fancyfoot[R]{}
\pagestyle{fancy}
\newcommand{\gt}{\textgreater}
\newcommand{\lt}{\textless}
\usepackage{listings}
\usepackage{xcolor}
\lstset{
    basicstyle=\ttfamily\small,
    keywordstyle=\color{blue},
    commentstyle=\color{gray},
    stringstyle=\color{red},
    numbers=left,
    numberstyle=\color[gray]{0.7}\ttfamily\small,
    stepnumber=1,
    numbersep=8pt,
    frame=single,
    showstringspaces=false,
    tabsize=4,
    breaklines=true
}
\usepackage{subcaption}

\begin{document}
	\begin{titlepage}
		\setlength{\parindent}{0ex}
		
		\begin{center}
			\textsc{
				\vspace{1ex}
                Научно-исследовательский университет ИТМО \\
				\vspace{0.5ex}
				Факультет систем управления и робототехники \\
				\vspace{0.5ex}
			}
		\end{center}
		
		\vspace{45mm}
		
		\begin{center}
			Отчет по лабораторной работе №6\\
			\textbf{Синтез оптимального управления \\
			$H_\infty$-оптимизация} \\[0.5ex]
			Вариант 21
		\end{center}
		
		\vspace{50mm}
		
		\begin{minipage}{.45\linewidth}
			Выполнил студент
            \\
			Проверил преподаватель
		\end{minipage}
		\hfill
		\begin{minipage}{.52\linewidth}
			\begin{flushright}
				Мовчан Игорь Евгеньевич, R3480
                \\
				Парамонов Алексей Владимирович
			\end{flushright}
		\end{minipage}
		
		\vfill
		\begin{center}
			Санкт-Петербург
			\\
			2026
		\end{center}
		
	\end{titlepage}

	\tableofcontents
	\clearpage
	
	\section{Синтез $H_\infty$-оптимального регулятора}
	Рассмотрим возмущённый линейный объект управления, описываемый уравнением состояния с начальными условиями $x(0)$:
	\[
		\dot{x} = Ax + Bu + B_f f, \quad x(0) = \begin{bmatrix}
			1 & 0
		\end{bmatrix}^T
	\]

	Задача заключается в построении H$_\infty$-оптимального регулятора по состоянию вида $u = Kx$. Расчет матрицы коэффициентов $K$ производится на основе решения алгебраического уравнения Риккати:
	\[
	\begin{cases}
		A^T P + PA + Q - PBB^T P + \gamma^{-2}PB_f B_f^T P = 0 \\
		K = -B^T P
	\end{cases}
	\]

	Необходимо экспериментально определить минимальное значение коэффициента подавления помех $\gamma = \gamma_{min}$, при котором существует положительно полуопределённое решение $P \succcurlyeq 0$.

	Матрицы системы и коэффициенты возьмём согласно варианту:
	\[
		A = \begin{bmatrix}
			3 & 1 \\
			1 & 1
		\end{bmatrix}, \quad
		B = \begin{bmatrix}
			1 \\
			2
		\end{bmatrix}, \quad
		B_f = \begin{bmatrix}
			3 \\
			4
		\end{bmatrix}, \quad
		Q = \begin{bmatrix}
			3 & 0 \\
			0 & 5
		\end{bmatrix}
	\]

	\section{Определение $\gamma_{min}$ и расчет регулятора}

	Минимальное значение $\gamma_{min}$, при котором существует стабилизирующее положительно полуопределенное решение $P$ уравнения Риккати, является ключевым параметром, определяющим максимально достижимый уровень подавления возмущений.

	С помощью нескольких попыток поиска было найдено искомое минимальное значение:
	\[
		\gamma_{min} \approx 3.5
	\]

	Для найденного $\gamma_{min}$ было решено уравнение Риккати, что дало следующую симметричную матрицу $P$:
	\[
		P = \begin{bmatrix}
			115.195 & -14.584 \\
			-14.584 &  3.696
		\end{bmatrix}
	\]

	На основе этой матрицы была вычислена матрица коэффициентов обратной связи $K$ для H$_\infty$-оптимального регулятора:
	\[
		K = -B^T P = \begin{bmatrix}
			-86.027 & 7.192
		\end{bmatrix}
	\]

	Полученный регулятор $u=Kx$ обеспечивает робастную стабилизацию системы при заданном уровне внешних возмущений.

	Моделирование при внешнем возмущении
	\[
		f(t) = 10 \sin 6t + 5 \cos 2t + 4 \cos 3t + 3 \cos 8t
	\]

	приведено на рисунках \ref{fig:u} и \ref{fig:x}. Можем видеть, что управление успешно стабилизирует состояния $x(t)$ системы - внешнее возмущение $f(t)$ не разносит систему. Регулятор робастен, успех!
	\begin{figure}[h]
		\centering
		\includegraphics[width=0.8\textwidth]{images/u.png}
		\caption{График управляющего воздействия при оптимальном регуляторе}
		\label{fig:u}
	\end{figure}
	\begin{figure}
		\centering
		\includegraphics[width=0.8\textwidth]{images/x.png}
		\caption{Графики переменных состояния объекта при оптимальном регуляторе}
		\label{fig:x}
	\end{figure}

	\section{Расчёт норм передаточных функций}
	Для анализа качества подавления возмущений, действующих на систему, необходимо рассчитать $H_\infty$-нормы передаточных функций от входа возмущения $f$ до различных выходов системы. В данном случае, в качестве выходов рассматриваются компоненты вектора состояния $x_1$ и $x_2$, а также сам вектор состояния $x$.

	Замкнутая регулятором система описывается уравнением:
	\[
		\dot{x} = (A + BK)x + B_f f
	\]
	
	Пусть далее $A_{M} = A + BK$ — матрица замкнутой системы.

	Передаточные функции от возмущения $f$ до компонент состояния $x_1$ и $x_2$ (обозначим их $G_1(s)$ и $G_2(s)$) имеют вид:
	\begin{align*}
		G_1(s) &= C_1 (sI - A_{M})^{-1} B_f \\
		G_2(s) &= C_2 (sI - A_{M})^{-1} B_f
	\end{align*}
	где $C_1 = \begin{bmatrix} 1 & 0 \end{bmatrix}$ и $C_2 = \begin{bmatrix} 0 & 1 \end{bmatrix}$.

	Передаточная функция от возмущения $f$ до всего вектора состояния $x$ (обозначим ее $G_x(s)$) тогда:
	\[
		G_x(s) = I(sI - A_{M})^{-1} B_f = (sI - A_{M})^{-1} B_f
	\]

	$H_\infty$-норма передаточной функции определяется как пиковое значение её АЧХ по всем частотам:
	\[
		\| G(s) \|_{\infty} = \sup_{\omega \in \mathbb{R}} \sigma_{\max}(G(j\omega))
	\]
	где $\sigma_{\max}(\cdot)$ обозначает максимальное сингулярное число. Для скалярных передаточных функций $G_1(s)$ и $G_2(s)$ это эквивалентно максимуму модуля $|G(j\omega)|$.

	В итоге были получены следующие значения норм:
	\begin{itemize}
		\item $H_\infty$-норма передаточной функции от $f$ до $x_1$:
		\[
			\| G_1(s) \|_{\infty} = \| C_1(sI - A_M)^{-1}B_f \|_{\infty} \approx 0.1079
		\]
		\item $H_\infty$-норма передаточной функции от $f$ до $x_2$:
		\[
			\| G_2(s) \|_{\infty} = \| C_2(sI - A_M)^{-1}B_f \|_{\infty} \approx 1.4602
		\]
		\item $H_\infty$-норма передаточной функции от $f$ до $x$:
		\[
			\| G_x(s) \|_{\infty} = \| (sI - A_M)^{-1}B_f \|_{\infty} \approx 1.4642
		\]
	\end{itemize}

	Отметим, что значение $H_\infty$-нормы представляет собой коэффициент усиления <<в худшем случае>> от входного возмущения $f(t)$ к выходному сигналу.
	\begin{itemize}
		\item Норма $\| G_1(s) \|_{\infty}$ показывает, во сколько раз в худшем случае амплитуда гармонического возмущения на некоторой частоте будет усилена в установившемся режиме для состояния $x_1$.
		\item Аналогично, $\| G_2(s) \|_{\infty}$ характеризует максимальное усиление для состояния $x_2$.
		\item Норма $\| G_x(s) \|_{\infty}$ даёт обобщенную оценку усиления для всего вектора состояния $x$.
	\end{itemize}

	Сравнивая значения $\| G_1(s) \|_{\infty}$ и $\| G_2(s) \|_{\infty}$, можно сделать вывод о том, что состояние $x_2$ является более чувствительной к влиянию возмущения $f(t)$. Это можно было видеть и на графиках.

	\section{Выводы}
	В ходе лабораторной работы был успешно выполнен синтез $H_\infty$-оптимального регулятора для линейной динамической системы. Было численно определено минимальное значение коэффициента подавления возмущений $\gamma_{min}$, после чего, решив соответствующее уравнение Риккати, была найдена матрица коэффициентов регулятора. Таким образом, была решена задача робастной стабилизации объекта управления по полному вектору состояния.

	
\end{document}